\subsection*{Riesgos operacionales}
\addcontentsline{toc}{subsection}{Riesgos operacionales}

La gestión de riesgos operacionales es un proceso cíclico continuo que busca identificar, evaluar y controlar los riesgos relacionados a operaciones diarias de la empresa, lo que conlleva a su aceptación, mitigación o elusión.Esta gestión implica la supervisión constante del riesgo operativo, incluyendo la posibilidad de pérdidas derivadas de fallas en los procesos y sistemas internos, factores humanos u otros eventos externos \parencite{Microbank01}.

\begin{enumerate}
	\item \textbf{Fallas técnicas y de vulnerabilidades:} Se refiere a la necesidad de que la infraestructura de la plataforma sea lo suficientemente robusta para evitar caídas del sistema o problemas de rendimiento, especialmente en momentos de alta demanda, así como la aparición de vulnerabilidades. Para ello, se implementará una infraestructura sólida, junto con protocolos de ciberseguridad y la realización de pruebas constantes.
	
	\item \textbf{Gestión del crecimiento:} A medida que se expandan las operaciones de la compañía, un crecimiento acelerado y no planificado podría generar dificultades en la gestión de la plataforma. Tanto la atención al cliente como la calidad del servicio podrían verse afectadas. Por esta razón, se propone llevar a cabo una expansión gradual, planificada y controlada de las operaciones hacia nuevos mercados.
	
	\item \textbf{Incumplimiento de objetivos:} Se refiere a la incapacidad de la empresa para cumplir con los objetivos y plazos establecidos en el desarrollo de nuevas iniciativas o proyectos. Para mitigar este riesgo, se realizará un estudio detallado de los proyectos que ingresen a la empresa, así como un análisis adecuado de los recursos humanos involucrados en su ejecución.
\end{enumerate}


